\section{训练能保证什么不能保证什么}
\label{sec:18-Synthesis/01-System-Lifecycle/04}

梯度检查通过了，损失还是不降。

一个排序模型的训练日志：交叉熵从 \(0.693\) 掉到 \(0.52\)，然后在四十个 epoch 里几乎不动。数值梯度与解析梯度逐元素比对，相对误差在 \(10^{-7}\) 量级，实现没有问题。有人把学习率调大十倍，损失直接变成 \texttt{nan}；调小十倍，四十个 epoch 只挪了 \(0.003\)。换一个随机种子重跑，最终损失是 \(0.47\)——比原来好了一大截，可谁也说不清为什么。

这类排查之所以耗时，是因为``训练不动''这一个症状底下压着三种性质完全不同的问题，而它们各自需要的动作互相冲突。把它们分开是这一步唯一重要的事：

\textbf{优化保证的是``在这个目标函数下找到了一个好的点''，它对目标函数写得对不对一无所知}；凸问题保证局部即全局，非凸问题连这个都没有；而\textbf{数值上收敛不等于数学上收敛}——步长、条件数、浮点精度都能让一条完全正确的算法给出错误的答案。

\begin{center}
\begin{tikzpicture}[x=1cm,y=1cm,font=\small,>=stealth]
  \draw[black!55] (0.4,2) -- (11.6,2);
  \foreach \x in {0.9,4.3,7.7,11.1} \fill[blue!60!black] (\x,2) circle (0.09);
  \node[align=center,font=\footnotesize,anchor=north,black!85] at (0.9,1.8) {训练脚本\\吐出的参数};
  \node[align=center,font=\footnotesize,anchor=north,black!85] at (4.3,1.8) {训练目标的\\真正最小};
  \node[align=center,font=\footnotesize,anchor=north,black!85] at (7.7,1.8) {真实风险的\\最小};
  \node[align=center,font=\footnotesize,anchor=north,black!85] at (11.1,1.8) {原始诉求的\\最优};
  \draw[<->,black!55] (1.1,2.4) -- (4.1,2.4);
  \draw[<->,black!55] (4.5,2.4) -- (7.5,2.4);
  \draw[<->,black!55] (7.9,2.4) -- (10.9,2.4);
  \node[align=center,font=\footnotesize,black!70] at (2.6,3.1) {步长、停止准则\\浮点、非凸地形};
  \node[align=center,font=\footnotesize,black!70] at (6.0,3.1) {有限样本\\（下一篇）};
  \node[align=center,font=\footnotesize,black!70] at (9.4,3.1) {代理指标\\（第一篇）};
  \draw[black!45,fill=blue!8] (1.1,-0.35) rectangle (4.1,0.35);
  \node[font=\footnotesize,black!80] at (2.6,0) {本篇只管这一段};
\end{tikzpicture}

\emph{从训练结果到真正想要的东西之间隔着三段缺口，性质各不相同。训练这一步只负责最左边那一段；把它做到极致，另外两段一寸都不会缩短。}
\end{center}

\subsection{优化对目标函数是否正确不发表意见}

上图最左段之外的一切，优化器都不知道。它拿到一个可微函数就去降低它，降得越低越尽职。如果损失里少写了一个负号、如果类别权重写反了、如果某个正则项的系数量纲错了三个数量级，优化器不会有任何异常表现——它会兢兢业业地把那个错误的目标优化到底，然后交出一条漂亮的下降曲线。

这件事听上去显然，但它导致的实际错误极多，因为下降曲线本身太有说服力了。损失在降，人就会认为方向是对的。一条便宜的防线是先在一小批数据上把模型训到过拟合：如果几百个样本都记不住，那基本可以断定实现或目标有问题，而不必再去调优化器。\hyperref[sec:08-Numerical-Analysis/07-Numerical-Methods-for-ML/02]{``自动微分给出导数，梯度检查建立信任''}给出的另一条防线是数值梯度比对，它能抓住实现错误，但抓不住``目标本身就写错了''——这两件事只有前者能被自动化检查覆盖。

顺带一提，梯度写错还有一类不报错的形态是形状对了但布局约定不一致。\hyperref[sec:10-Matrix-Calculus/01-Derivative-Conventions/02]{``形状先于公式，布局必须声明''}讲的正是这个：分子布局与分母布局差一个转置，在方阵情形下维度检查照样通过，训练也照样跑，只是每一步走的方向不对。

\subsection{凸与非凸之间隔着一条真实的分界}

\hyperref[sec:09-Convex-Optimization/01-Optimization-Foundations/02]{``局部更好为何还不够''}提出的是整个优化理论的出发问题：算法只能看见当前点附近，凭什么它找到的局部最好就是全局最好。凸性给出的回答干脆——\hyperref[sec:09-Convex-Optimization/03-Convex-Functions/02]{``切平面为何是全局下界''}说明凸函数在任意一点的切平面都落在函数图像下方，于是梯度为零这个纯局部的条件，直接推出该点在整个定义域上最小。这是一条定理，不是经验规律，换来它的条件是函数确实凸。

非凸问题里这条推理链断在第一环。梯度为零只说明到了驻点，而驻点可以是局部极小、局部极大或鞍点。\hyperref[sec:09-Convex-Optimization/08-Nonconvex-and-Stochastic-Optimization/02]{``二阶结构区分局部极小与鞍点''}指出在高维空间里鞍点比局部极小常见得多——一个驻点要成为局部极小，需要 Hessian 的所有特征值同号，而维数越高这件事越难同时发生。训练平台期很多时候不是陷进了局部极小，是在一片近乎平坦的区域里方向不明。

有意思的是，非凸并不意味着一切保证消失。\hyperref[sec:09-Convex-Optimization/08-Nonconvex-and-Stochastic-Optimization/03]{``PL 条件让部分非凸问题恢复全局速率''}给出的是一类中间地带：只要梯度的大小能控制住当前值与最小值之差，就算函数不凸，也能证明线性收敛。深度网络的训练地形在某些区域被认为具有类似性质，但这属于对现象的建模而非对实际网络的证明，不应当当成保证来用。

\subsection{数值上收敛与数学上收敛是两件事}

回到开头那个案例。梯度是对的，目标也没写错，损失却卡住——最常见的原因是条件数。

\hyperref[sec:08-Numerical-Analysis/01-Floating-Point-and-Error/03]{``条件数衡量问题本身的敏感性''}讲的是问题固有的敏感度：输入的微小扰动会被放大多少倍。在优化里它的作用是另一面——Hessian 的条件数决定了等高线有多扁，而梯度方向在扁长的谷里几乎垂直于谷底方向。\hyperref[sec:09-Convex-Optimization/06-Optimization-Algorithms/01]{``梯度下降真正难在步长''}把这件事量化成一条判断：步长的上限由最大曲率决定，收敛速度由最小曲率决定，两者之比就是条件数，于是条件数为 \(10^4\) 时每次迭代的进展会慢上四个数量级。这解释了为什么``调大就炸、调小就停''，那不是玄学，是同一个数从两头卡住了可用区间。而特征标准化之所以有用，正是因为它直接改善这个数——这又一次说明上一篇的表示选择会渗透到下游。

浮点是另一条独立的失效通道。\hyperref[sec:08-Numerical-Analysis/01-Floating-Point-and-Error/04]{``前向误差、后向误差与稳定算法''}提供的是判断算法好坏的正确标准：一个稳定算法给出的解，是某个被轻微扰动过的问题的精确解。这句话把责任分成两半——问题本身病态（条件数大）与算法不稳定是两件事，前者换算法救不了，后者换算法能救。实践里 \texttt{float32} 换 \texttt{float64} 之后结果明显变化，说明累积误差已经和真实信号可比，此时该查的是有没有做减法抵消、有没有直接求逆而不是解方程、概率有没有在对数域计算。\hyperref[sec:08-Numerical-Analysis/07-Numerical-Methods-for-ML/01]{``在对数域中稳定计算概率''}讲的那套变换是这一类问题最常见的解法。

随机性带来第三种模糊。小批量梯度是全量梯度的无偏估计，但方差不为零，于是即使真实梯度已经很小，随机梯度仍然把参数推来推去。\hyperref[sec:09-Convex-Optimization/08-Nonconvex-and-Stochastic-Optimization/04]{``SGD 在期望中下降但噪声会留下平台''}给出的判断是：固定步长下损失会降到一个由噪声水平决定的平台就不再下降，想继续降必须让步长趋于零。这意味着看到平台期时，先该问的是``这是噪声地板还是真的收敛了''，而降一次学习率就能区分——如果损失应声再降一截，那之前的平台是噪声造的。

\hyperref[sec:14-Deep-Learning/03-Optimization-and-Training/01]{``训练动力学与数值稳定''}把上面这些在深度网络里的具体形态汇到一处，其中批量统计量的漂移值得单独记住：归一化层在训练时用批内统计量、推理时用滑动平均，两者不一致会让一个训练良好的模型在推理时表现异常，而这个现象在训练日志上完全看不见。

\subsection{这一步做错时，故障长什么样}

损失曲线出现长平台，这是最模糊的症状，因为它至少对应三种原因：条件数过大（降学习率无效、做预处理有效）、随机噪声地板（降学习率立刻有效）、真的进了平坦区（两者都无效）。区分它们的成本很低，而不区分就调参的成本很高。

换一个初始化种子结果差很多，说明训练落在了地形上非常不同的位置，或者优化过程本身方差过大。这件事本身不是错误，但它意味着单次运行的结果不足以支撑任何比较——两个方案的差距如果小于种子之间的差距，那个比较就没有内容。这是下一篇的话题在训练阶段的预演。

\texttt{float32} 换 \texttt{float64} 结果改变，或者同一份代码在不同硬件上给出不同结果，说明计算已经处在数值边缘。此时任何基于该结果的结论都要打折，包括那些看起来无关的结论——因为误差是沿计算图传播的，它不会只影响一个数。

还有一种沉默的形态：训练一切正常，损失降到很低，模型却在推理路径上表现异常。这通常不是优化的问题，是训练图与推理图不一致（归一化统计量、dropout 的开关、数据管线里的随机增强）。它的特征是离线复现训练时的输入能得到正确输出，走推理接口就不行。

\hyperref[sec:09-Convex-Optimization/08-Nonconvex-and-Stochastic-Optimization/06]{``优化误差、统计误差与泛化不能混为一谈''}把这一篇与下一篇的分界写成了一句可用的判断：损失降不下去是优化问题，训练低而测试高是统计问题。这两句话对应完全不同的动作，而它们唯一的共同点是都表现为``模型不好用''。

训练结束，参数已经定下来了。接下来要做的事是给出一个数，并且用这个数向别人承诺这个系统在没见过的数据上会怎样——那是一句关于未来的断言，它的成立条件比训练严格得多。
